\begin{helper}
In the activation--priority output map of
\noderef{RandomIndependentSetSampling}, a vertex $v$ belongs to the output
set
\[
I(A,\pi)=\{u\in A:\pi(w)<\pi(u)\text{ for every activated neighbor }w
\text{ of }u\}
\]
if and only if $v$ is activated and every activated neighbor of $v$ has
priority smaller than the priority of $v$.
\end{helper}
\begin{proof}
The output set in \noderef{RandomIndependentSetSampling} is defined by
filtering the full finite vertex set with the predicate
\[
u\in A\quad\text{and}\quad
\forall w,\; w\in A\Rightarrow uw\in E(G)\Rightarrow \pi(w)<\pi(u).
\]
Membership in a filtered finite set is equivalent to membership in the
ambient finite set together with satisfaction of the filtering predicate.
The ambient finite set is the full vertex set, so $v$ automatically belongs
to it.  Substituting $u=v$ in the displayed predicate gives exactly the two
conditions: $v\in A$, and every activated neighbor $w$ of $v$ has
$\pi(w)<\pi(v)$.  Conversely, if those two conditions hold, then $v$
satisfies the filtering predicate and hence belongs to the output set.
\end{proof}
